\section{Fixed-target cohomological image filtrations}

\begin{lemma}[The finite original tail Rees module]
\label{lem:fc-formal-functions-fixed-tail}
\leanok
\lean{EnvelopeProperty.FormalFunctionsFixedTarget.tailSpan, EnvelopeProperty.FormalFunctionsFixedTarget.polynomialMap, EnvelopeProperty.FormalFunctionsFixedTarget.polynomialMap_lof_of_le, EnvelopeProperty.FormalFunctionsFixedTarget.polynomialMap_lof, EnvelopeProperty.FormalFunctionsFixedTarget.polynomialMap_tail_smul, EnvelopeProperty.FormalFunctionsFixedTarget.tailReesMap, EnvelopeProperty.FormalFunctionsFixedTarget.tailSpan_finite}
\uses{lem:fc-rees-cech-homogeneous-action,lem:fc-rees-cech-degree-transitions,lem:fc-rees-cech-homology-module}
Fix $n,q\geq0$.  In the actual Rees module
$M=\bigoplus_d H^q(I^d\mathcal O_X)$, let $T_n$ be the Rees submodule
generated by the original summands of degrees $n+d$.  The original
transition map sends a homogeneous element $y\in H^q(I^{n+d}\mathcal O_X)$
to its image in $H^q(I^n\mathcal O_X)$ times $t^d$.  This extends to a
Rees-linear map on $T_n$.  If the original scheme is projective over a
Noetherian base, $T_n$ is finite over the original Rees algebra.
\end{lemma}
\begin{proof}
First define an $R$-linear map on the whole direct sum by the stated
formula in degrees at least $n$ and by zero in lower degrees.  For
$a\in I^e$ and a summand of degree $n+d$, genuine multiplication has
degree $e+n+d$, and lowering its image to degree $n$ gives $a$ times
the original transition image.  Its polynomial degree is $e+d$, as
required.  Additive monomial induction proves this scalar equality for
every Rees scalar on each original summand of degree at least $n$.
For any additive map between modules, the elements on which it
commutes with every scalar form a submodule: additivity proves closure
under sums, and for a scalar $b$ the scalar equality on $bx$ follows
by applying the equalities on $x$ for $ab$ and $b$.  This submodule
contains the generators of $T_n$, so restriction to $T_n$ is
Rees-linear.  Projective Rees cohomology makes $M$ finite, and the
Rees algebra of an ideal in a Noetherian ring is Noetherian.  Its
submodule $T_n$ is consequently finite.
\end{proof}

\begin{lemma}[Stability at every fixed original cohomology degree]
\label{lem:fc-formal-functions-fixed-target-stable}
\leanok
\lean{EnvelopeProperty.FormalFunctionsFixedTarget.imageFiltration, EnvelopeProperty.FormalFunctionsFixedTarget.imageFiltration_zero, EnvelopeProperty.FormalFunctionsFixedTarget.polynomialMap_mem, EnvelopeProperty.FormalFunctionsFixedTarget.tailReesMap_range, EnvelopeProperty.FormalFunctionsFixedTarget.imageFiltration_stable}
\uses{lem:fc-formal-functions-fixed-tail,lem:fc-formal-functions-degree-finite,lem:fc-formal-functions-rees-finite-image}
For a projective scheme over a Noetherian ring, every fixed-target
filtration
$F^n_d=\operatorname{im}(H^q(I^{n+d}\mathcal O_X)\to H^q(I^n\mathcal O_X))$
is a stable $I$-filtration of the same original finite module
$H^q(I^n\mathcal O_X)$, beginning at the whole module.
\end{lemma}
\begin{proof}
Transition composition proves that the images decrease with $d$.
The multiplication and lowering identity gives
$IF^n_d\subseteq F^n_{d+1}$, and the identity transition gives
$F^n_0=H^q(I^n\mathcal O_X)$.  The tail Rees map has image precisely
the associated submodule $\bigoplus_d F^n_d$: coefficients of all
images lie in the prescribed transition images, and each homogeneous
polynomial with a coefficient in $F^n_d$ comes from its original
degree-$(n+d)$ preimage.  Since these homogeneous polynomials generate
the associated module, the map is surjective onto that entire module.
The tail is finite, so this associated module is finite.  Original
projective cohomology makes the target finite over the Noetherian base,
and the associated-module criterion now proves stability.
\end{proof}
